% Generated draft: requires mathematical and source review.
\documentclass[11pt]{article}
\usepackage[margin=25mm]{geometry}
\usepackage{fontspec}
\setmainfont{Latin Modern Roman}
\usepackage{amsmath,amssymb,mathtools,unicode-math}
\setmathfont{Latin Modern Math}
\usepackage{enumitem,needspace,xurl,hyperref}
\setlength{\parindent}{0pt}
\setlength{\parskip}{5pt}
\setlength{\emergencystretch}{3em}
\newcommand{\N}{\mathbb{N}}
\newcommand{\Z}{\mathbb{Z}}
\newcommand{\Q}{\mathbb{Q}}
\newcommand{\R}{\mathbb{R}}
\newcommand{\C}{\mathbb{C}}
\newcommand{\F}{\mathbb{F}}
\newcommand{\A}{\mathbb{A}}
\newcommand{\PP}{\mathbb{P}}
\newcommand{\E}{\mathbb{E}}
\newcommand{\eps}{\varepsilon}
\newcommand{\dd}{\,\mathrm{d}}
\begin{document}
\section*{UnsolvedMath ID: 800064}
\textbf{Source problem code:} catalogue entry
\subsection{Definitions and mathematical statement}
For a smooth three-dimensional manifold \(M\), the usual meaning of \emph{irreducible} is that every smoothly embedded sphere \(S^2\subset\operatorname{int}(M)\) bounds a smoothly embedded closed three-ball in \(\operatorname{int}(M)\).

The supplied question asks: if \((\xi_-,\xi_+)\) is a \emph{tight positive contact pair} on \(M\), in the sense and under the hypotheses of the preceding question in the source list, must \(M\) be irreducible? Under the smooth three-dimensional interpretation above, the requested conclusion is that every such embedded sphere bounds such a three-ball.

\textbf{Unresolved qualification.} The preceding question is absent from the supplied record. Consequently, the intended class of manifolds, the definition of a tight positive contact pair, and the hypotheses implicit in that reference cannot be specified from this record alone. In particular, no compactness, boundary, orientability, connectedness, regularity, or compatibility assumptions are added here. The question remains conditional on recovering those source definitions and hypotheses.

As background, the supplied source reports that refinements of the Eliashberg--Thurston theorem by Colin and Bowden imply that any Reebless aspherical \(C^0\)-foliation approximates to such pairs. This is retained as a reported contextual claim; the record does not specify the approximation topology or the precise meaning of the foliation qualifications. It does not explicitly require the pair in the question to arise from this approximation construction.
\subsection{Short English statement}
The question, attributed in the supplied list to Colin and Firmo, asks whether a manifold carrying a tight positive contact pair must be irreducible. The definition of the pair and the ambient hypotheses depend on a preceding question that was not supplied, so this is an incomplete statement-only draft requiring mathematical review.
\subsection{Sources}
\begin{itemize}
\item Supplied local UnsolvedMath registry record, permanent ID \(800064\), exact source problem code \texttt{catalogue entry}, titled \emph{Foliations and Diffeomorphism Groups --- Question 5}. The registry attributes the record to Nariman and Yazdi, with proposed year \(2025\), and labels it \texttt{open}; this is a supplied status label, not an independently verified assessment.
\item Nariman and Yazdi, \emph{Problem List on Foliations and Diffeomorphism Groups} (2025), Question 5, \url{https://arxiv.org/abs/2503.18274}. The supplied provenance records access on 2026-10-16 and extraction by \texttt{arxiv\_latex}; these are supplied provenance details, not a claim of independent access or verification.
\item The source attributes the question to Colin and Firmo and cites the unresolved reference key \texttt{colin2011paires}. No full bibliographic entry or direct link for that reference is supplied.
\item The contextual approximation claim is attributed in the supplied statement to refinements of the Eliashberg--Thurston theorem by Colin and Bowden. No separate bibliographic references for these refinements are supplied.
\end{itemize}
\end{document}
